\chapter{Conclusion}
\label{chap:conc}

This project demonstrates that an appropriately executed 2D Attention U-Net architecture, when trained and validated using a rigorous per-volume procedure, can deliver competitive abdominal CT segmentation of multiple organs using just the computing power of a free-tier cloud GPU. The final model achieves a mean Dice score of 0.8769 and HD95 of 7.59~mm across 13 abdominal organs on the FLARE22 benchmark, accomplishing this in under 80 minutes of training on a single T4 GPU.

The design of this project was influenced by key decisions. Using attention gates on all skip connections made a big difference for small organs. It stopped the decoder from getting distracted by things that were not important in the background. The way we combined cross-entropy and Dice loss helped us from the start of training. It made sure we had gradients and it also helped with the problem of class imbalance in abdominal CT scans. We used cosine annealing to decrease the learning rate. This allowed the learning rate to decay smoothly, enabling continued fine-tuning through the full 60 epochs without overshooting. We also did something called Test-Time Augmentation with a four-way flip. This added approximately $+$0.06 Dice at zero additional training cost. We made a change to how we evaluated the results. We used to look at each slice one by one but this was not accurate because some slices were empty. So we started looking at each volume as a whole. This gave us a more realistic idea of how well the segmentation was working. The corrected per-volume evaluation protocol --- replacing the na\"\i{}ve per-slice average that inflated scores through trivially empty slices --- ensured that every number in this report reflects genuine segmentation quality.

It is helpful to compare with nnU-Net~3D. On large, well-defined organs the 2D model comes within 0.01 Dice of the 3D baseline. On 7 out of 13 organs, our 2D model numerically exceeds the published nnU-Net~3D baseline --- a result we attribute largely to the 2D model's ability to process full-resolution in-plane slices without the patch-level downsampling that 3D methods require to fit GPU memory. While the comparison gives an indication, it is not controlled --- nnU-Net figures come from a different evaluation protocol. The pancreas and duodenum are more challenging cases due to their irregular shapes and limited tissue contrast. Volumetric models can provide the inter-slice context that these structures require natively.

There are three important lessons. Firstly, the rigour of an evaluation is as important as the architecture. For example, fixing the per-slice Dice bug changed the reported score from an inflated $\sim$0.88 to an honest 0.8164 (no TTA) --- a correction that is honest but uncomfortable. The second point is that data volume is the dominating factor at this scale --- scaling from 30 to 50 training volumes added more Dice than any single architectural or scheduling change. Third, the gap between 2D and 3D is not fixed; it gets significantly narrower when sufficient data, correct evaluation, and modern training recipes are provided to 2D methods.

In summary, the Attention U-Net trained in this project is a practical, accurate and efficient solution for multi-organ abdominal CT segmentation in the context of real hardware limitations.